\chapter{The Integers}

\section{Cyclic Sets}

Natural numbers count elements of finite sets. The empty set has zero elements, the set of solutions to the equation $3x-3=0$ has exactly one element, the booleans have exactly two elements, true and false, and so on.

Natural numbers have been around for at least as long as humans count. But what do negative numbers count? Surely, there are no sets with $-5$ elements! One explanation of negative numbers is that they represent deficits. In other words, integers can be thought of as \emph{differences} between natural numbers. This point of view is completely valid, and we will come back to it. There is, however, a more structural and conceptually compelling way of thinking about integers, using cyclic sets. Cyclic sets are an often recurring structure in mathematics. Some of the prettiest proofs of \zcrefnamed{theorem:fermats-little-theorem} and \zcrefnamed{theorem:wilsons-theorem} use them.

\begin{figure}[b]
  \centering
  \begin{tikzpicture}[
    alt={
      A set A with 12 elements labelled a sub 1 through a sub 12 is displayed with exactly one outgoing arrow from each element, indicating a mapping. The mapping is a bijection, so that each element also has exactly one incoming arrow. The mapping sends a sub 1 to a sub 5 to a sub 9 to a sub 1; a sub 2 to itself; a sub 3 to a sub 4 to a sub 8 to a sub 12 to a sub 7 to a sub 3; and a sub 6 to a sub 11 to a sub 10 to a sub 6. Notice that there are four disjoint cycles in this mapping, which are called orbits.
      },
    element/.style={circle, draw, thick, minimum size=6.5mm, inner sep=0pt, font=\small},
    set/.style={ellipse, draw=CornflowerBlue, thick, minimum height=35mm, minimum width=80mm},
    arrow/.style={->, thick, shorten >=1mm, shorten <=1mm}
    ]

    \node[set] (A) at (0,0) {};

    \node[element] (a1)  at ($(A.center)+(-3.0,  0.6)$) {$a_1$};
    \node[element] (a2)  at ($(A.center)+(-0.3,  1.1)$) {$a_2$};
    \node[element] (a3)  at ($(A.center)+( 1.0,  1.2)$) {$a_3$};
    \node[element] (a4)  at ($(A.center)+( 2.4,  0.9)$) {$a_4$};

    \node[element] (a5)  at ($(A.center)+(-1.6,  0.5)$) {$a_5$};
    \node[element] (a6)  at ($(A.center)+(-0.7,  0.0)$) {$a_6$};
    \node[element] (a7)  at ($(A.center)+( 1.8,  0.2)$) {$a_7$};
    \node[element] (a8)  at ($(A.center)+( 3.2, -0.1)$) {$a_8$};

    \node[element] (a9)  at ($(A.center)+(-2.5, -0.6)$) {$a_9$};
    \node[element] (a10) at ($(A.center)+(-1.2, -1.1)$) {$a_{10}$};
    \node[element] (a11) at ($(A.center)+( 0.5, -0.9)$) {$a_{11}$};
    \node[element] (a12) at ($(A.center)+( 1.6, -1)$) {$a_{12}$};

    \draw[arrow] (a1) to[bend left=10] (a5);
    \draw[arrow] (a5) to[bend left=10] (a9);
    \draw[arrow] (a9) to[bend left=10] (a1);

    \draw[arrow] (a2) to[out=340,in=280,looseness=5] (a2);

    \draw[arrow] (a3) to[bend left=10] (a4);
    \draw[arrow] (a4) to[bend left=10] (a8);
    \draw[arrow] (a8) to[bend left=10] (a12);
    \draw[arrow] (a12) to[bend left=10] (a7);
    \draw[arrow] (a7) to[bend left=10] (a3);

    \draw[arrow] (a6) to[bend left=10] (a11);
    \draw[arrow] (a11) to[bend left=10] (a10);
    \draw[arrow] (a10) to[bend left=10] (a6);

  \end{tikzpicture}  
  \caption[A bijection on a $12$-element set]{A bijection on a $12$-element set.}
  \label{figure:bijection-12-element-set}
\end{figure}

When we map out a bijection $f:A\to A$ from a finite set $A$ to itself by displaying an arrow between each element and its value under $f$, as we did in \zcref{figure:bijection-12-element-set}, we discover that the finite set $A$ gets partitioned into cycles. The iterated action of $f$ on an element $a$ is mapped out by tracing along the arrows, which eventually loops back to itself. An intuitive reason for this phenomenon is that by iterating $f$ often enough, we obtain a sequence
\begin{equation*}
  a,~f(a),~f^2(a),~\ldots
\end{equation*}
which must eventually have repeated values since the set $A$ was assumed to be finite. If $f^n(a)$ is the first $n$ such that $f^n(a)=f^m(a)$ for a later value of $m$, then we must have $n=0$. Indeed, if $n>0$ then we find that
\begin{equation*}
  f(f^{n-1}(a))=f(f^{m-1}(a)),
\end{equation*}
which implies by the assumption that $f$ is a bijection that $f^{n-1}(a)=f^{m-1}(a)$. This is impossible by the minimality of $n$, so we must have $n=0$. We conclude that the iterations of $f$ trace out cycles. Such cycles are called \emph{orbits}.

\begin{definition}
  Consider a bijection $f:X\cong X$. We say that $x,y\in X$ are \define{in the same orbit} of $f$ if there exists a natural number $k$ such that $f^k(x)=y$ or $f^k(y)=x$.

  A subset $U\subseteq X$ is an \define{orbit} of $f$ if there exists an element $x\in X$, such that $U$ consists of exactly the elements in the same orbit as $x$.
\end{definition}

The bijection on the $12$-element set in \zcref{figure:bijection-12-element-set}, for example, has four such orbits. Since bijections on a set partition the set into cyclic orbits, we can now define a \emph{cyclic set} to be a set equipped with a bijection that has exactly one orbit, so that all the elements of the set fit in the same cycle. In \zcref{figure:cyclically-ordered-set} we displayed an example of a cyclic set.

\begin{figure}[t]
  \centering
  \begin{tikzpicture}[
    alt={
      A set A with 12 elements labeled 0 through 11, and arrows between them indicating a mapping. This mapping has a single cycle: It maps 0 to 1 to 2 to 3, and so on until it reaches 11, which is sent back to 0.
      },
    thick, every node/.style={circle, draw, minimum size=7mm, inner sep=1pt}]
    \def\n{12}
    \def\radius{2.5cm}
    
    \foreach \i in {1,...,\n} {
      \pgfmathsetmacro{\value}{int(\i-1)}
      \node (N\i) at ({90-360/\n*(\i-1)}:\radius) {$\value$};
    }
    
    \foreach \i in {1,...,\n} {
      \pgfmathsetmacro{\angle}{90-360/\n*(\i-1)}
      \draw[->] 
      (\angle-10:\radius)
      arc[start angle=\angle-10, delta angle=-360/\n+20, radius=\radius];
    }
\end{tikzpicture}  
  \caption[A cyclic set with $12$ elements]{\label{figure:cyclically-ordered-set}A cyclic set with $12$ elements.}
\end{figure}

\begin{definition}
  A \define{cyclic set} is a pair $(X,s)$ consisting of a set $X$ and a bijection $s:X\cong X$ that has exactly one orbit.
\end{definition}

\begin{remark}
  Since orbits are assumed to be inhabited---that is, nonempty---subsets of a set $X$, it follows that cyclic sets are always inhabited. Indeed, if $U$ is the unique orbit of a cyclic set, and $U$ is inhabited by an element $x\in U$, then we have $x\in X$. Thus, empty sets are not considered to be cyclically structured.
\end{remark}

\begin{example}\label{example:standard-finite-cyclic-sets}
  Every inhabited standard finite set can be given the structure of a cyclic set. The cyclic structure on the standard finite set $[n+1]$ is given by
  \begin{equation*}
    0\mapsto 1 \mapsto 2 \mapsto \cdots \mapsto n \mapsto 0.
  \end{equation*}
  To define this function precisely, observe that there are two functions $\iota,\sigma:[n]\to[n+1]$ given by
  \begin{equation*}
    \iota(x)=x\qquad\text{and}\qquad \sigma(x)=x+1.
  \end{equation*}
  These allow us to define the successor function $S$ and its inverse $T$ as follows:
  \begin{equation*}
    S(x)=
    \begin{cases}
      \sigma(y) & \text{if $x = \iota(y)$,} \\
      0    & \text{if $x = n$,}
    \end{cases}
    \qquad\qquad
    T(x)=
    \begin{cases}
      \iota(y) & \text{if $x = \sigma(y)$,} \\
      n & \text{if $x = 0$.}
    \end{cases}
  \end{equation*}
  Thus, we obtain a cyclic set $([n+1],0,S)$ for each $n$.
\end{example}

\section{Maps Preserving Cyclic Structure}

\begin{figure}[t]
  \centering
  \begin{tikzpicture}[
    alt={
      Two cyclic sets are displayed, with a mapping between them. The first cyclic set is displayed as an outer cycle, around the second cyclic set, so as to indicate the winding of the first cyclic set around the second cyclic set. The first cyclic set has 14 elements labeled 0 through 13 and the second cyclic set has 7 elements labeled 0 through 6. The mapping therefore winds around twice, sending 0 to 3, 1 to 4, and so on until the assignment 13 to 2. 
      },
    every node/.style={
      circle,draw,fill=white,inner sep=1pt,minimum size=6mm,thick}
    ]
    \newdimen\radiusInner 
    \newdimen\radiusMiddle
    \newdimen\radiusOuter
        
    \def\radiusInner{1.5cm}
    \def\radiusMiddle{3.5cm}
    \def\radiusOuter{4.3cm}
    
    \foreach \i in {0,...,6} {
      \pgfmathsetmacro{\ang}{-360/7*\i}
      \pgfmathtruncatemacro{\highlabel}{7+\i}
      \pgfmathtruncatemacro{\targetlabel}{mod(3+\i,7)}
      \node (N1\i) at (\ang:\radiusInner)  {$\targetlabel$};
      \node (N2\i) at (\ang:\radiusMiddle) {$\highlabel$};
      \node (N3\i) at (\ang:\radiusOuter)  {$\i$};
    }

    \foreach \i in {0,...,6} {
      \pgfmathsetmacro{\angle}{-360/7*(\i-1)}
      \draw[->,thick,shorten <=2pt,shorten >=2pt] 
      (\angle-12:\radiusInner)
      arc[start angle=\angle-12, delta angle=-360/7+24, radius=\radiusInner];
    }

    \foreach \i in {0,...,6} {
      \draw [|->,thick,shorten >=6pt,shorten <=6pt] (N2\i) -- (N1\i);
    }

    \foreach \i in {0,...,5} {
      \pgfmathsetmacro{\angle}{-360/7*\i}
      \draw[->,thick,shorten <=2pt,shorten >=2pt] 
      (\angle-6:\radiusMiddle)
      arc[start angle=\angle-6, delta angle=-360/7+12, radius=\radiusMiddle];
    }

    \foreach \i in {0,...,5} {
      \pgfmathsetmacro{\angle}{-360/7*\i}
      \draw[->,thick,shorten <=2pt,shorten >=2pt] 
      (\angle-5:\radiusOuter)
      arc[start angle=\angle-5, delta angle=-360/7+10, radius=\radiusOuter];
    }

    \draw[->,thick,shorten >=3pt,shorten <=3pt] (N36) to[out=360/7-95,in=96] (N20);
    \draw[-,line width=3pt,white,shorten >=3pt,shorten <=3pt] (N26) to[out= 360/7-96,in=95] (N30);
    \draw[->,thick,shorten >=2pt,shorten <=2pt] (N26) to[out= 360/7-96,in=95] (N30);
  \end{tikzpicture}
  \caption[A cyclic-structure map from a $14$-element set to a $7$-element set]{A cyclic-structure map from a $14$-element set to a $7$-element set.}
  \label{figure:example-cyclic-structure-map}
\end{figure}

An important organizational principle of structural mathematics is that when we define a certain kind of structured sets, such as cyclic sets in our case, then it is good practice to also define the \emph{structure-preserving maps} between them. Doing so allows us to learn more about how such structured sets relate to each other. Thus, we will now introduce cyclic-structure maps. As we will see shortly, maps preserving the cyclic structure on a set have some interesting properties, and they feature prominently in our structural definition of the set of integers.

\begin{definition}
  Consider two cyclic sets $(X,s)$ and $(Y,t)$. A map $f:X\to Y$ is said to be a \define{cyclic-structure map} if
  \begin{equation*}
    f(s(x))=t(f(x))
  \end{equation*}
  for every $x\in X$. We will write $\mathrm{hom}((X,s),(Y,t))$ for the set of all cyclic-structure maps from $(X,s)$ to $(Y,t)$.
\end{definition}

For example, if $(X,s)$ is a cyclic set, then $s$ is always a cyclic-structure map onto itself. To see this, we simply note that the equality
\begin{equation*}
  s(s(x))=s(s(x))
\end{equation*}
holds trivially. In \zcref{figure:example-cyclic-structure-map} we displayed an example of a cyclic-structure map from a $14$-element set to a $7$-element set. More generally, the idea of a cyclic-structure map is that we can overlay the elements of a cyclic set $X$ onto the elements of a cyclic set $Y$ so that the cyclic structure of the overlay matches the cyclic structure of $Y$. This means that in a cyclic-structure map, we can't skip any elements.

\begin{figure}[t]
  \centering
  \begin{tikzpicture}[
    alt={
      Two abstractions of cyclic sets are displayed, one above the other. The upper one maps down to the second one, and winds around five times.
      },
    declare function={f(\x)=0.2*sin(\x)+\x/900;},
    line cap=butt,
    line join=miter,
    rubout/.style={%
      /utils/exec=\tikzset{rubout/.cd,#1},
      decoration={
        show path construction,
        curveto code={
          \draw [white,line width=\pgfkeysvalueof{/tikz/rubout/line width}+2*\pgfkeysvalueof{/tikz/rubout/halo}] 
          (\tikzinputsegmentfirst) .. controls
          (\tikzinputsegmentsupporta) and (\tikzinputsegmentsupportb)  ..(\tikzinputsegmentlast); 
          \draw [thick,shorten <=-0.1pt,shorten >=-0.1pt] (\tikzinputsegmentfirst) .. controls
          (\tikzinputsegmentsupporta) and (\tikzinputsegmentsupportb) ..(\tikzinputsegmentlast);  
        }
      }
    },
    rubout/.cd,
    line width/.initial=2pt,
    halo/.initial=0.5pt
    ]
    
    \draw[rubout={line width=1.5pt,halo=0.5pt},decorate]
    
    plot[variable=\x,domain=-50:1690,samples=200,smooth] ({cos(\x)},{f(\x)}) to[out=0,in=195] cycle;
    \draw[thick] (0,-2) arc(-90:270:1cm and 0.2cm);
    \draw[thick,->]  (0,-0.4) -- (0,-1.4);
  \end{tikzpicture}  
  \caption[Cyclic-structure maps]{Cyclic-structure maps are defined by ``winding'' the source cyclic set around the target cyclic set.}
  \label{figure:cyclic-structure-maps}
\end{figure}

\begin{theorem}
  Consider two cyclic sets $(X,s)$ and $(Y,t)$ with a cyclic-structure map $f:X\to Y$ between them. If $X$ is an $m$-element set and $Y$ is an $n$-element set, then $m$ is a multiple of $n$.
\end{theorem}

\begin{proof}
  For each $y\in Y$ consider the set
  \begin{equation*}
    f^{-1}(y) := \{x\in X\mid f(x)=y\}
  \end{equation*}
  consisting of all the elements in $X$ that have the value $y$ under $f$. Then the bijection $s:X\to X$ restricts to a bijection
  \begin{equation*}
    s : f^{-1}(y) \to f^{-1}(s(y)).
  \end{equation*}
  Thus it follows that all the sets of the form $f^{-1}(y)$ have the same size. Since $X$ is the union of these sets, it follows that the size of $X$ is $n$ times the size of the set $f^{-1}(y)$.
\end{proof}

Cyclic-structure maps have the special property that if they share an equal value for a given input, then they are equal everywhere. Indeed, suppose that $f,g:X\to Y$ are two cyclic-structure maps between cyclic sets $(X,s)$ and $(Y,t)$ such that $f(x)=g(x)$ for some $x\in X$. By the cyclic structure of $X$, it follows that for every $x'\in X$, there exists a natural number $k$ such that $s^k(x)=x'$ or $s^k(x')=x$. Thus there are two cases to consider. If $s^k(x)=x'$ then it follows that
\begin{equation*}
  f(x')=f(s^k(x))=s^k(f(x))=s^k(g(x))=g(s^k(x))=g(x').
\end{equation*}
On the other hand, if $s^k(x')=x$ then it follows that $f(x')=g(x')$ if and only if $s^k(f(x'))=s^k(g(x'))$, which follows from the calculation
\begin{equation*}
  s^k(f(x'))=f(s^k(x'))=f(x)=g(x)=g(s^k(x'))=s^k(g(x')).
\end{equation*}
The previous observation can be used to prove the following theorem:

\begin{theorem}
  Consider two cyclic sets $(X,s)$ and $(Y,t)$. If the set
  \begin{equation*}
    \mathrm{hom}((X,s),(Y,t))
  \end{equation*}
  of cyclic-structure maps from $(X,s)$ to $(Y,t)$ is inhabited, then it is itself a cyclic set.
\end{theorem}

\begin{proof}
  The function $h\mapsto t\circ h$ is a bijection from $\hom((X,s),(Y,t))$ to itself. To see that it has exactly one orbit, consider two cyclic-structure maps $f,g:X\to Y$, and let $x_0\in X$. Then there exists an integer $k$ such that
  \begin{equation*}
    t^k(f(x_0))=g(x_0)\qquad\text{or}\qquad t^k(g(x_0))=f(x_0).
  \end{equation*}
  In the first case it follows that $t^k\circ f=g$, while in the second case it follows that $t^k\circ g=f$. Thus it follows that $f$ and $g$ are in the same orbit of the map $h\mapsto t\circ h$.
\end{proof}

The organizational principle of defining structured sets along with their structure-preserving maps comes from \emph{category theory}, a mathematical discipline which emphasizes structural mathematics. What we have defined here is the category of cyclic sets. In this course we won't rely on category theory, but the curious reader might find Emily Riehl's book \emph{Category Theory in Context} of interest \cite{riehl2016}. Riehl is a world-famous professor in our department.

\section{A Structural Definition of the Integers}

\begin{figure}[t]
  \centering
  \begin{tikzpicture}[
    alt={
      The infinite cyclic set has for each integer an element, and the cyclic structure maps each integer to its successor. The infinite cyclic set thus extends infinitely in both directions, the negative direction and the positive direction.
    },
    scale=.9,
    thick]
    \def\n{6}       
    \def\xsep{.65in}
    
    \foreach \i in {0,...,\n} {
      \pgfmathtruncatemacro{\value}{\i-3}
      \node[circle, draw, minimum size=7mm, inner sep=1pt,font=\small] (N\i) at (\i*\xsep,0) {$\value$};
    }
    
    \pgfmathtruncatemacro{\nMinusOne}{\n-1}
    
    \foreach \i in {0,...,\nMinusOne} {
      \pgfmathtruncatemacro{\next}{\i+1}
      \draw[->, thick, shorten >=1mm, shorten <=1mm]
      (N\i) -- (N\next);
    }
    
    \draw[->, thick, dashed, shorten >=1mm, shorten <=4mm]
    ($(N0)+(-\xsep,0)$) -- (N0);
    
    \draw[->, thick, dashed, shorten >=4mm, shorten <=1mm]
    (N\n) -- ++(\xsep,0);
  \end{tikzpicture}
  \caption[The infinite cyclic set of integers]{The infinite cyclic set of integers.}
  \label{figure:integers}
\end{figure}

Although the examples we have discussed so far are finite, the definition of cyclic sets isn't restricted to finite sets. We could indeed imagine a very large cyclic set that extends infinitely in both directions, as in \zcref{figure:integers}. If such a set has a designated zero element, we could use it as the starting point to describe the other elements as the iterated successors and the iterated predecessors. This way, we would expect to find one element for each integer. Alternatively, as we are about to do now, we could take the infinite cyclic set as the definition of the set of integers!

\begin{definition}\label{definition:integers}
  The set of \define{integers} is specified to be a set $\mathbb{Z}$ equipped with a distinguished \define{zero} element $0\in \mathbb{Z}$ and a \define{successor} bijection $S:\mathbb{Z}\cong\mathbb{Z}$ such that for every set $X$, every element $x_0$ and every bijection $h:X\cong X$ uniquely determine a function $f:\mathbb{Z}\to X$ such that
  \begin{enumerate}
  \item $f(0)=x_0$, and
  \item $f(S(a))=h(f(a))$ for every $a\in\mathbb{Z}$.
  \end{enumerate}
  Such functions are said to be defined by \define{integer recursion}. We will commonly write $a\mapsto a+1$ for the successor function $S$, and $a\mapsto a-1$ for its inverse.
\end{definition}

For example, integer recursion can be used to define \emph{iterated bijections} $f^k:A\cong A$ for any bijection $f:A\cong A$. The idea is entirely analogous to the idea of iterated functions, where a function $f:A\to A$ could be iterated any natural number of times. For iterated bijections, we restrict to the case of a bijection $f:A\cong A$, which allows us to also consider the \emph{negative iterations} of $f$ by using the inverse of $f$. In order to formally define iterated bijections, we need the following theorem:

\begin{theorem}
  For any bijection $f: A\cong B$ and any set $X$, there are bijections
  \begin{align*}
    \blank\circ f & : (B\cong X)\cong (A\cong X) \\
    f\circ\blank & : (X\cong A)\cong (X\cong B).
  \end{align*}
\end{theorem}

\begin{proof}
  Suppose $f:A\cong B$ is a bijection with inverse $f^{-1}$. Then for any bijection $g:B\cong X$ we have
  \begin{equation*}
    g\circ f\circ f^{-1}=g,
  \end{equation*}
  and for any $h:A\cong B$ we have
  \begin{equation*}
    h\circ f^{-1}\circ f=h.
  \end{equation*}
  This shows that the function $\blank\circ f^{-1}$ is inverse to the function $\blank\circ f$. Similarly $f^{-1}\circ\blank$ is inverse to the function $f\circ\blank$.
\end{proof}

\begin{definition}
  Consider a set $X$ equipped with a bijection $f:X\cong X$. Then we define for each integer $k$ the \define{$k$\supth{} iteration} of $f$ by
  \begin{align*}
    f^{0} & := \mathrm{id} \\
    f^{k+1} & := f\circ f^k.
  \end{align*}
  This definition is possible, since composing with $f$ on the left is a bijection $(X\cong X)\cong(X\cong X)$.
\end{definition}

The specification of the set of integers with integer recursion also allows us to formulate an induction principle for the integers. This induction principle states what one should do in order to show that a proposition $P(x)$ is true for all integers $x$. The induction principle for the integers was first formulated by Angelo Margaris \cite{margaris1961}.

\begin{theorem}[Integer Induction Principle]\label{theorem:integer-induction}
  Consider a property $P(a)$ parametrized by an integer $a$. In order to prove that $P(a)$ is true for all $a$, it suffices to prove:
  \begin{enumerate}
  \item The property $P(0)$ is true.
  \item The logical equivalence $P(a)\Leftrightarrow P(a+1)$ holds for all $a\in\mathbb{Z}$. 
  \end{enumerate}
\end{theorem}

Note the similarity with the induction principle of the natural numbers: We have to prove a base case and an inductive step. However, in the inductive step we have to prove two directions:
\begin{equation*}
  P(a) \Rightarrow P(a+1)\qquad\text{and}\qquad P(a+1)\Rightarrow P(a).
\end{equation*}
In other words, there are two steps in the inductive step: the \emph{forward step} and the \emph{backward step}. In the forward step, you will assume $P(a)$ as the inductive hypothesis as usual, and the task is to prove that $P(a+1)$ holds. In the backward step, you will assume $P(a+1)$ is true, and the task is to prove that $P(a)$ holds.
In \zcref{section:integer-arithmetic} we will use the integer induction principle to prove the laws for integer arithmetic.

\begin{proof}
  Consider a property $P(a)$ of the integers, such that $P(0)$ is true and for which the logical equivalence
  \begin{equation*}
    P(a)\Leftrightarrow P(a+1)
  \end{equation*}
  holds. In order to show that $P(a)$ is true for every integer $a$, we will use the recursion principle for $\mathbb{Z}$. Thus, we have to construct a set $A$ equipped with an element $a_0\in A$ and a bijection $s:A\cong A$. We define the set $A$ by
  \begin{equation*}
    A := \{a\in\mathbb{Z}\mid P(a)\},
  \end{equation*}
  so that $A$ consists of exactly those integers $a$ for which $P(a)$ holds. Clearly, we have $0\in A$, so we choose $a_0:=0$. In order to define a function $s:A\to A$, observe that for $a\in A$, the property $P(a)$ holds by assumption, which implies $P(a+1)$ so that $a+1\in A$. Thus, we simply define $s(a):=a+1$.

  To see that $s$ is a bijection, we define an inverse $t$. For any $a\in A$, observe that the logical equivalence $P(a-1)\Leftrightarrow P(a)$ holds, since $(a-1)+1=a$, so it follows that $P(a-1)$ holds. This allows us to define $t(a):=a-1$. It follows immediately from the definition of $t$ that the identities $s(t(a))=a$ and $t(s(a))=a$ hold. Thus we see that $t$ is an inverse of $s$, which concludes the proof that $s$ is a bijection on $A$.

  By the specification of $\mathbb{Z}$, it follows that there is a unique function $f:\mathbb{Z}\to A$ such that
  \begin{equation*}
    f(0)=0\qquad\text{and}\qquad f(S(a))=S(f(a)).
  \end{equation*}
  Now we claim that $f(a)=a$ for all $a\in\mathbb{Z}$. To see this, note that by the inclusion of $A\subseteq \mathbb{Z}$ we may regard $f$ as a function $\mathbb{Z}\to\mathbb{Z}$ satisfying the equations
  \begin{equation*}
    f(0)=0\qquad\text{and}\qquad f(a+1)=s(f(a))=f(a)+1.
  \end{equation*}
  However, there is a known function satisfying these equations: the \emph{identity function} $x\mapsto x$, which is often denoted by $\mathrm{id}$. Since $f$ is uniquely determined by its equations, which are also satisfied by the identity function, it follows that $f$ and the identity function must be the same. Thus it follows that $f(a)=a$ for all $a\in\mathbb{Z}$. Furthermore, since we have $f(a)\in A$ by definition, we obtain that $a\in A$ for every integer $a$. Thus we conclude that the property $P(a)$ holds for all $a\in\mathbb{Z}$.
\end{proof}

The induction principle for the integers may be used, for instance, to prove that the integers form a cyclic set.

\begin{theorem}
  The pair $(\mathbb{Z},S)$ consisting of the set $\mathbb{Z}$ of integers and the successor bijection $S$ is a cyclic set. Furthermore, for any cyclic set $(Y,t)$ with a distinguished element $y_0$, there is a unique cyclic-structure preserving map $f$ from $(\mathbb{Z},S)$ to $(Y,t)$ satisfying $f(0)=y_0$.
\end{theorem}

\begin{proof}
  We show by integer induction that every integer is in the orbit of $0$. Clearly $0$ is in its own orbit, so the base case is trivial. For the inductive step, we have to show that $x$ is in the orbit of $0$ if and only if $S(x)$ is.

  Suppose first that $x$ is in the orbit of $0$. If $S^k(0)=x$ for some natural number $k$, then $S^{k+1}(0)=S(x)$. Otherwise, $S^k(x)=0$ for some natural number $k$. The case $k=0$ gives $x=0$ and is already covered. We may therefore write $k=k'+1$, and we obtain $S^{k'}(S(x))=0$. In either case, $S(x)$ is in the orbit of $0$.

  Conversely, suppose that $S(x)$ is in the orbit of $0$. If $S^k(0)=S(x)$ for some natural number $k$, we distinguish two cases. If $k=0$, then $S(x)=0$, so $x$ is in the orbit of $0$. If $k=k'+1$, injectivity of $S$ gives $S^{k'}(0)=x$. Alternatively, if $S^k(S(x))=0$ for some natural number $k$, then $S^{k+1}(x)=0$. Thus $x$ is in the orbit of $0$, completing the induction.

  The final claim follows from the integer recursion principle.
\end{proof}

\begin{theorem}
  Let the inclusion of the naturals into the integers $i:\mathbb{N}\to\mathbb{Z}$ be defined by
  \begin{align*}
    i(0) & := 0, \\
    i(n+1) & := S(i(n)).
  \end{align*}
  Then $0$ is not of the form $S(i(n))$, for any natural number $n$.
\end{theorem}

\begin{proof}
  Suppose that $0=S(i(n))=i(n+1)$ for some natural number $n$. Recall from \zcref{example:standard-finite-cyclic-sets} that the set $\{0,\ldots,n+1\}$ has the structure of a cyclic set with successor map $S$. Integer recursion therefore gives us a function $f:\mathbb{Z}\to\{0,\ldots,n+1\}$ satisfying the equations
  \begin{equation*}
    f(0)=0 \qquad\text{and}\qquad f(a+1)=S(f(a)).
  \end{equation*}
  Since $S(x)=x+1$ for each $0\leq x\leq n$, it follows that $f(i(x))=x$ for each $0\leq x\leq n+1$. Now we obtain a contradiction with \zcref{theorem:is-not-successor-zero-nat}:
  \begin{equation*}
    0=f(0)=f(i(n+1))=n+1. \qedhere
  \end{equation*}
\end{proof}

\section{Constructing the Integers from the Natural Numbers}

We have given a structural definition of the set of integers as the universal cyclic set. We think of this as the \emph{specification} of the integers. In this section we will construct a set that meets this specification.

\begin{figure}[t]
  \centering
  \begin{tikzpicture}[
    alt={
      The lattice of points m comma n is displayed in a two dimensional coordinate system, where m and n range over the natural numbers. The points on the axis of abscissas as well as the points on the axis of ordinates are marked as solid points, while the points in the positive quadrant, where both coordinates are positive integers, are marked as open points. The set of lattice points thus defined mark the integers, according to the Cartesian definition of the integers.
      },
    scale=1,
    line join=round,
    line cap=round
    ]
    
    \draw[->] (-0.3,0) -- (10.5,0) node[right] {$x$};
    \draw[->] (0,-0.3) -- (0,4.5) node[above] {$y$};
    
    \foreach \x in {0,1,2,3,4,5,6,7,8,9,10} {
      \foreach \y in {0,1,2,3,4} {
        \fill[CornflowerBlue] (\x,\y) circle (2pt);
      }
    }

    \foreach \x in {1,2,3,4,5,6,7,8,9,10} {
      \foreach \y in {1,2,3,4} {
        \draw[CornflowerBlue,fill=white,thick] (\x,\y) circle (2pt);
      }
    }    
    
  \end{tikzpicture}
  \caption[The set of integers along the axis of the positive quadrant]{The set $Z$ consists of the lattice points on the coordinate axes whose coordinates are natural numbers. Along the axis of abscissas we find the nonnegative integers in their usual positions, while along the axis of ordinates we find the nonpositive integers.
}
  \label{figure:implementation-integers}
\end{figure}
  
\begin{definition}
  We define the set $\mathbb{Z}_C$ of \define{cartesian integers} by
  \begin{equation*}
    \mathbb{Z}_C:=\{(a,b)\in\mathbb{N}\times\mathbb{N}\mid ab=0\}.
  \end{equation*}
  The \define{zero-element} of $\mathbb{Z}_C$ is defined to be the element $(0,0)$. We define the \define{successor function} $s:\mathbb{Z}_C\to \mathbb{Z}_C$ recursively on $b$ by
  \begin{align*}
    s(a,0) & := (a+1,0) \\
    s(a,b+1) & := (a,b).
  \end{align*}
\end{definition}

\begin{theorem}
  The successor function $s:\mathbb{Z}_C\to \mathbb{Z}_C$ is a bijection.
\end{theorem}

\begin{proof}
  We prove that it is a bijection by constructing its inverse $t:\mathbb{Z}_C\to \mathbb{Z}_C$. The function $t$ is defined recursively on $a$ by
  \begin{align*}
    t(0,b) & := (0,b+1) \\
    t(a+1,b) & := (a,b).
  \end{align*}
  The proof that $s\circ t=\mathrm{id}$ is by induction on $a$. We have $s(t(0,b))=s(0,b+1)=(0,b)$ proving the base case, and we have
  \begin{equation*}
    s(t(a+1,b))=s(a,b)=(a+1,b)
  \end{equation*}
  since $b=0$ in this case.

  Similarly, the proof that $t\circ s=\mathrm{id}$ is by induction on $b$. We have $t(s(a,0))=t(a+1,0)=(a,0)$ proving the base case, and we have
  \begin{equation*}
    t(s(a,b+1))=t(a,b)=(a,b+1)
  \end{equation*}
  since $a=0$ in this case. This completes the proof that $t$ is an inverse of $s$, and hence $s$ is a bijection.
\end{proof}

\begin{theorem}
  The set $\mathbb{Z}_C$ equipped with the distinguished zero element $(0,0)$ and the successor function $s$ satisfies the specification of the integers in \zcref{definition:integers}.
\end{theorem}

\begin{proof}
  Consider a set $X$ equipped with a point $x_0\in X$ and a bijection $h:X\cong X$. Then we define a function $f:\mathbb{Z}_C\to X$ by recursion on $a$ and $b$:
  \begin{align*}
    f(0,0) & := x_0 \\
    f(a+1,0) & := h(f(a,0)) \\
    f(0,b+1) & := h^{-1}(f(0,b)).
  \end{align*}
  This function satisfies $f(0,0)=x_0$ by definition. We claim that
  \begin{equation*}
    f(s(a,b))=h(f(a,b))
  \end{equation*}
  for every $(a,b)\in \mathbb{Z}_C$. The proof is by induction on $b$. For the base case, observe that the equalities $f(s(a,0))=f(a+1,0)=h(f(a,0))$ hold by definition. For the inductive step, note that the equality $f(0,b+1)=h^{-1}(f(0,b))$ implies that $h(f(0,b+1))=f(0,b)$. Using this, we obtain
  \begin{equation*}
    f(s(0,b+1))=f(0,b)=h(f(0,b+1)).
  \end{equation*}
  Thus far, we have proven that there exists a function $f:\mathbb{Z}_C\to X$ satisfying
  \begin{equation*}
    f(0,0)=x_0\qquad\text{and}\qquad f(s(a,b))=h(f(a,b)).
  \end{equation*}
  It remains to show that this function $f$ is unique.
\end{proof}

Knowing that there is a set that meets the specification in \zcref{definition:integers}, we shall now abandon the set $\mathbb{Z}_C$ again, and we will consider $\mathbb{Z}$ to be an arbitrary set as specified in \zcref{definition:integers}.

\section{Integer Arithmetic}\label{section:integer-arithmetic}

We can now define addition of two integers, just as we did for the natural numbers, by using the iterated successor bijection. Furthermore, since iterated bijections are again bijections, we obtain something that we didn't get for the natural numbers: For every integer $y$ there is a unique integer $x$ such that $x+a=y$. In other words, we can define differences of integers.

\begin{definition}
  We will define \define{addition} of two integers and the \define{additive inverse} of an integer:
  \begin{enumerate}
  \item For any two integers $a$ and $b$, we define
  \begin{equation*}
    a+b :=S^b(a),
  \end{equation*}
  where $S^b$ is the $b$\supth{} iteration of the successor function $S:\mathbb{Z}\cong\mathbb{Z}$. Henceforth, we will write $a+1$ for $S(a)$.
\item For any integer $a$, we define $-a$ to be the unique element in the set
  \begin{equation*}
    \{x\in\mathbb{Z}\mid x+a=0\}.
  \end{equation*}
  Such an element is indeed uniquely determined, since $S^a$ is a bijection.
  \end{enumerate}
\end{definition}

Before we start proving the laws of addition and the additive inverse of the integers, we will show how addition interacts with iteration.

\begin{theorem}
  For any bijection $f:X\cong X$ with inverse $f^{-1}$, and for any $x\in X$, we have
  \begin{align*}
    f^{a+b}(x) & =f^b(f^a(x)) \\
    f^{-a}(x) & = (f^{-1})^a(x).
  \end{align*}
\end{theorem}

\begin{proof}
  For the first identity, we will use \zcrefnamed{theorem:integer-induction} on $b$. It is clear that
  \begin{equation*}
    f^{a+0}(x)=f^a(x)=f^0(f^a(x)),
  \end{equation*}
  so the base case holds. To see that
  \begin{equation*}
    f^{a+b}(x)=f^b(f^a(x)) \qquad\Leftrightarrow\qquad f^{a+(b+1)}(x)=f^{b+1}(f^a(x))
  \end{equation*}
  note that $f^{a+(b+1)}(x)=f(f^{a+b}(x))$ and $f^{b+1}(f^a(x))=f(f^b(f^a(x)))$. Furthermore, since $f$ is a bijection, it follows that $f(y)=f(z)$ if and only if $y=z$. In our case we get:
  \begin{align*}
    f^{a+b}(x)=f^b(f^a(x)) & \Leftrightarrow f(f^{a+b}(x))=f(f^b(f^a(x))) \\
                           & \Leftrightarrow f^{a+(b+1)}(x)=f^{b+1}(f^a(x)).
  \end{align*}
  This completes the proof of the first identity. The second identity is proven by integer induction on $a$. Since $-0=0$, it follows that
  \begin{equation*}
    f^{-0}(x) = f^0(x)=x=(f^{-1})^0(x).
  \end{equation*}
  For the inductive step, we calculate in similar fashion:
  \begin{align*}
    f^{-a}(x)=(f^{-1})^a(x) & \Leftrightarrow f^{-(a+1)+1}(x) = f(f^{-1}((f^{-1})^a(x))) \\
    & \Leftrightarrow f(f^{-(a+1)}(x)) = f((f^{-1})^{a+1}(x)) \\
    & \Leftrightarrow f^{-(a+1)}=(f^{-1})^{a+1}(x). \qedhere
  \end{align*}
\end{proof}

\begin{theorem}
  Addition of integers satisfies the following laws:
  \begin{align*}
    (a+b)+c & = a+(b+c), & a+(b+1) & = (a+b)+1, \\
    a+0 & = a, & (a+1)+b & = (a+b)+1, \\
    0+a & = a, & a+b & = b+a. \\
    a-a & = 0,
  \end{align*}
\end{theorem}

\begin{proof}
  By the previous theorem, we calculate:
  \begin{equation*}
    (a+b)+c= S^c(S^b(a))=S^{b+c}(a)=a+(b+c).
  \end{equation*}
  Next, note that the unit law $a+0=a$ holds because
  \begin{equation*}
    a+0=S^0(a)=a
  \end{equation*}
  Similarly, we note that $0+a=a$ by induction on $a$. The base case is covered in the previous step. For the inductive step, we have to show that
  \begin{equation*}
    0+a=a\qquad\Leftrightarrow\qquad 0+(a+1)=a+1.
  \end{equation*}
  This can be seen by chaining the injectivity of the successor function and associativity:
  \begin{align*}
    0+a=a & \Leftrightarrow (0+a)+1=a+1 \\
          & \Leftrightarrow 0+(a+1)=a+1.
  \end{align*}
  The fact that $a-a=0$ follows from the observation that
  \begin{equation*}
    a-a=(0+a)-a=S^{-a}(S^a(0))=0.
  \end{equation*}
  The right successor law $a+(b+1)=(a+b)+1$ is an instance of the associativity law, which we have already shown. On the other hand, the left successor law $(a+1)+b=(a+b)+1$ is proven by induction on $b$.
  
  The two unit laws $a+0=a$ and $0+a=a$ combined settle the base case for the inductive proof that $a+b=b+a$. For the inductive step, we have to show that
  \begin{equation*}
    a+b=b+a \qquad\Leftrightarrow\qquad a+(b+1)=(b+1)+a.
  \end{equation*}
  This follows from the injectivity of the successor function and the left and right successor laws that we have previously shown.
\end{proof}

Next, we define multiplication of integers.

\begin{definition}
  We define the \define{multiplication operation} $a,b\mapsto ab$ on the integers by
  \begin{align*}
    a0 & := 0 \\
    a(b+1) & := ab+a
  \end{align*}
  This well-defined, since for any integer $a$, the function $x\mapsto x+a$ is a bijection.
\end{definition}

\begin{theorem}
  Multiplication of integers distributes from both sides over addition: For any three integers $a$, $b$, and $c$ we have
  \begin{equation*}
    a(b+c)=ab+ac\qquad\text{and}\qquad (a+b)c=ac+bc.
  \end{equation*}
\end{theorem}

\begin{proof}
  We will prove only the first equation, leaving the other equation as an exercise. To show that multiplication of the integers distributes from the left over addition, we will use integer induction on $c$ to show that for any two integers $a$ and $b$ we have $a(b+c)=ab+ac$.

  In the base case, we simply note that
  \begin{equation*}
    a(b+0)=ab=ab+0=ab+a0.
  \end{equation*}
  For the inductive step, we have to show that
  \begin{equation*}
    \forall_{a,b\in\mathbb{Z}}a(b+c)=ab+ac\qquad\Leftrightarrow \forall_{a,b\in\mathbb{Z}} a(b+(c+1))=ab+a(c+1).
  \end{equation*}
  For the forward implication, assume that $a(b+c)=ab+ac$ is true for any two integers $a$ and $b$. Then we have
  \begin{equation*}
    a(b+(c+1))=a((b+1)+c)=a(b+1)+ac=ab+a+ac=ab+a(c+1).
  \end{equation*}
  For the reverse implication, assume that $a(b+(c+1))=ab+a(c+1)$ is true for any two integers $a$ and $b$. Then we have
  \begin{equation*}
    a(b+c)=a((b-1)+(c+1))=a(b-1)+a(c+1)=ab-a+ac+a=ab+ac.\qedhere
  \end{equation*}
\end{proof}

\exercises

\begin{routine}
  \exercise Show that multiplication of integers is associative: For any three integers $a$, $b$, and $c$ we have
  \begin{equation*}
    (ab)c=a(bc).
  \end{equation*}
  \exercise Show that multiplication of integers is commutative: For any two integers $a$ and $b$ we have
  \begin{equation*}
    ab=ba.
  \end{equation*}
  \exercise Show that multiplication with a nonzero integer $a$ satisfies the cancellation law
  \begin{equation*}
    ax=ay\qquad\Rightarrow\qquad x=y.
  \end{equation*}
  \exercise Prove that for any two integers $a$ and $b$, either $b\leq a$ or $a\leq b$.
\end{routine}

\begin{foundational}
  \exercise Consider a bijection $f:X\cong X$, and let $x\in X$. Show that a subset $U\subseteq X$ is the orbit of $x$ if and only if $U$ is the intersection of all subsets $V\subseteq X$ containing $x$, such that $f(V)=V$.
  \exercise Consider a bijection $f:X\cong X$, and let $x\in X$. Show that the image of the unique map $h:\mathbb{Z}\to X$ determined by $x$ and $f$ is a cyclic set.
\end{foundational}

%%% Local Variables:
%%% mode: latex
%%% TeX-master: "number-theory"
%%% TeX-engine: luatex
%%% End:

